\section{The fibre axioms with enough intrinsic fibrations}

\begin{claim}{C.1}{Proved}{Universal weakness and fibrewise invariance}
Suppose $\mathcal M$ has pullbacks, $W$ contains isomorphisms and
satisfies two-out-of-three, and every map factors as a map in $W$
followed by an intrinsic fibration. Then
\[
 W\cap\operatorname{Fib}_W=W^u.
\]
Moreover every $W$-map between intrinsic fibrations over $S$ remains
weak after arbitrary change of base. Thus XII axiom (b) follows.
\end{claim}
\begin{proof}
Let $u:X\to Y$ be weak over $S$, with both structural maps intrinsic,
and factor an arbitrary $T\to S$ as $T\xrightarrow i T_0
\xrightarrow p S$, with $i\in W$ and $p$ intrinsic. Pullback of $u$
over $T_0$ is weak by the definition of $p$. The comparison maps from
pullback over $T$ to pullback over $T_0$ are weak because they are
pullbacks of $i$ along base changes of $X\to S$ and $Y\to S$.
The commutative square and two-out-of-three show that the pullback of
$u$ over $T$ is weak. This proves (b).

Apply the assertion to the arrow $q\to1_S$ in the slice when
$q\in W\cap\operatorname{Fib}_W$; it says precisely that $q$ is
universally weak. Conversely, if $q:X\to S$ is universally weak,
then for a weak arrow $T'\to T$ over $S$ both vertical arrows in the
square with upper edge $X\times_ST'\to X\times_ST$ are weak.
Two-out-of-three makes that upper edge weak, proving intrinsic
membership; the identity base change gives $q\in W$.
\end{proof}

\begin{claim}{C.2}{Proved}{Lead 3 without functorial factorization}
Under C.1's hypotheses, for every full subcategory
\[
 \mathcal F_S:=\operatorname{Fib}_W S\ \subseteq\
 \mathcal H\ \subseteq\ \mathcal M/S
\]
the inclusions induce equivalences on ordinary localizations and on
$\infty$-categorical localizations at the restricted $W_S$.
Consequently factorization alone implies all three axioms (a), (b), (c).
No global terminal object, functorial factorization, retract closure,
or strong saturation of $W$ is required.
\end{claim}
\begin{proof}
We give all hypotheses of the localization theorem being used.
For an ordinary category $\mathcal A$, the nerve is an
$\infty$-category with weak equivalences and fibrations in the sense
of \citet[Definitions~7.4.6 and 7.4.12]{RESCisinski2019} if it has a
terminal object, the weak class contains identities and has
two-out-of-three, and a class $F$ containing identities and closed
under composition has the following properties. Pullbacks of
$F$-maps between fibrant objects along maps from fibrant objects
exist and stay in $F$; the same pullbacks preserve $F\cap W$;
every map with fibrant codomain factors as $W$ followed by $F$.
Here fibrant means that the terminal map belongs to $F$.
Cisinski's Theorem~7.5.18 then states that inclusion of the full
subcategory of fibrant objects induces an equivalence of
$\infty$-categorical localizations.

Take $\mathcal A=\mathcal H$, and let $F$ consist of its morphisms
whose underlying map in $\mathcal M$ is intrinsic. The object
$1_S$ is terminal in $\mathcal H$ and belongs to $\mathcal F_S$.
Identities and compositions are in $F$, by the corresponding
unconditional properties of intrinsic fibrations. The fibrant
objects of $\mathcal H$ are exactly $\mathcal F_S$.
Suppose $X\to Y$ is in $F$, with $X,Y$ fibrant over $S$, and
$Z\to Y$ has $Z$ fibrant. Its pullback $P\to Z$ is intrinsic,
so $P\to S$ is intrinsic by composition. Thus $P$ belongs to
$\mathcal F_S$, hence to $\mathcal H$. Since $\mathcal H$ is full,
this is also a pullback in $\mathcal H$. If $X\to Y$ is weak,
C.1 makes its pullback weak. Finally, for a map $B\to A$ in
$\mathcal H$ with $A$ fibrant, factor its underlying map as
$B\xrightarrow w A'\xrightarrow q A$, where $w\in W$ and $q$
is intrinsic. Composition with $A\to S$ makes $A'$ fibrant,
so the entire factorization lies in $\mathcal H$.

All of Cisinski's hypotheses have now been verified within
$\mathcal H$. Importantly, his definition does not require arbitrary
pullbacks in $\mathcal H$. The required pullbacks have fibrant
vertices and have just been shown to remain in the smaller
subcategory. The theorem gives $L(\mathcal F_S)\simeq L(\mathcal H)$.
Apply the same argument to $\mathcal M/S$ to obtain the composite
equivalence $L(\mathcal F_S)\simeq L(\mathcal M/S)$; the second
inclusion is consequently an equivalence as well. Taking homotopy
categories gives the ordinary Gabriel--Zisman localizations.
C.1 identifies $W_S$ with $W_S^u$ on $\mathcal F_S$, giving
axiom (a) with exactly its stated weak class.
\end{proof}

\begin{claim}{C.3}{Refuted}{Five objects suffice without factorization}
Axioms (a) and (b), even at every base and for strongly saturated $W$,
do not imply (c) or the intrinsic factorization hypothesis. A
counterexample has five objects and sixteen arrows. Thus the
six-object poset example is not minimal among finite categories.
\end{claim}
\begin{proof}
Let $\mathcal M$ have objects $0,A,B,C,S$. The object $0$ is initial,
with its unique arrows to the other objects. Besides these arrows
and the five identities, there are
\[
 \alpha,\beta:A\rightrightarrows B,\quad \gamma:A\to C,
 \quad b:B\to S,\quad c:C\to S,\quad x,y:A\rightrightarrows S,
\]
with the only nonidentity composition relations not involving $0$
given by $b\alpha=c\gamma=x$ and $b\beta=y$. Put
\[
 W=\{1_0,1_A,1_B,1_C,1_S,\alpha,\beta,\gamma\}.
\]
These formulas define the complete composition law; associativity
involves only identities or arrows from $0$ and follows directly.
The functor to the three-object chain $0<m<s$ that sends $A,B,C$
to $m$ and $S$ to $s$ has $W$ as the inverse image of identities.
It proves two-out-of-three, and also strong saturation: any arrow
inverted in $\mathcal M[W^{-1}]$ must map to an identity in this
chain, hence must already belong to $W$.

Each slice is a poset. Over $S$ its six objects are
$0_S,x,y,b,c,1_S$, with Hasse relations
\[
 0_S<x,y,\qquad x<b,c,\qquad y<b,\qquad b,c<1_S.
\]
The slice over $B$ is the diamond $0_B<\alpha,\beta<1_B$;
over $C$ it is the chain $0_C<\gamma<1_C$; over $A$ it is
$0_A<1_A$; over $0$ it is terminal. Each has all binary meets,
so every cospan of $\mathcal M$ has a pullback. In particular
$\alpha$ and $\beta$ pull back to $0$, while $b$ and $c$ pull
back to $A$ with projections $\alpha,\gamma$.

The intrinsic fibrant objects in the slices are exactly
\[
\begin{array}{c|l|l}
\text{base}&\text{slice objects}&\text{intrinsic fibrant objects}\\\hline
0&1_0&1_0\\
A&0_A,1_A&0_A,1_A\\
B&0_B,\alpha,\beta,1_B&0_B,1_B\\
C&0_C,\gamma,1_C&0_C,\gamma,1_C\\
S&0_S,x,y,b,c,1_S&0_S,b,1_S
\end{array}
\]
To check the only nontrivial row, pull back the three weak arrows
$x\to b$, $y\to b$, $x\to c$ in the slice over $S$.
The object $b$ preserves all three, producing $\alpha,\beta,1_A$.
The candidates $x,y,c$ fail, respectively, on $y\to b$,
$x\to b$, and $y\to b$, producing $0\to A$ in each case.
The bottom object and the identity always pass. The other rows
follow from the same diamond and chain tests. These tests include
all possible structure maps to the base, not only weak arrows
whose codomain is the base.

On the fibrant subcategory the weak arrows are only identities,
except at base $C$, where $\gamma$ is also weak. Its pullbacks
over $C$ are $1_0,1_A,\gamma$. This proves (b) at all bases.
For (a), only $B$ and $S$ need replacement. Over $B$ send
$\alpha,\beta$ to $1_B$ and fix $0_B,1_B$; the identity has a
weak natural transformation to this retraction. Over $S$ define
the two endofunctors by
\[
\begin{array}{c|rrrrrr}
z&0_S&x&y&b&c&1_S\\\hline
Rz&0_S&b&b&b&b&1_S\\
Tz&0_S&x&y&b&x&1_S .
\end{array}
\]
Both preserve the order and the weak arrows. There are natural
transformations $1\leftarrow T\to R$ with weak components, and
$R,T$ restrict to the identity on $\{0_S,b,1_S\}$. They exhibit
the localized inclusion of the fibrant subcategory as an
equivalence, at both the ordinary and the enhanced level.

The full intermediate subcategory on $\{0_S,b,c,1_S\}$ is a
diamond and has only identity weak arrows. Its fibrant subcategory
is the chain $\{0_S,b,1_S\}$, also with only identity weak arrows.
The inclusion fails essential surjectivity at $c$, so (c) fails.
Finally the only weak arrow with source $C$ is $1_C$, whereas
$c:C\to S$ is not intrinsic. It has no required factorization.
The program verifies the category laws, all sixty-six cospans and
their universal properties, every intrinsic test, every test in
(b), and both displayed deformation certificates. Its complete
output is \texttt{code/BC/verification\_output.json}.
\end{proof}

\begin{claim}{C.4}{Proved}{Why an object bound is not a finite category search bound}
There are infinitely many inequivalent finite categories with
pullbacks and one object. A finite terminating exhaustive search
through six objects therefore needs an additional arrow bound or
a structural reduction. Every finite category with pullbacks has
only monomorphisms, but this does not bound its number of arrows.
\end{claim}
\begin{proof}
The one-object categories of cyclic groups $C_m$, for arbitrary
positive integers $m$, have pullbacks: a cospan of invertible
arrows has a pullback obtained using either inverse. They are
pairwise inequivalent. To prove the structural assertion, every
slice of a finite category with pullbacks is finite and has
terminal object and binary products. For any objects $X,Y$ in
such a slice, the product powers $Y^n$ give
$|\operatorname{Hom}(X,Y^n)|=|\operatorname{Hom}(X,Y)|^n$.
There is a common finite bound on all hom-set cardinalities,
so $|\operatorname{Hom}(X,Y)|\le1$. Each slice is therefore a
preorder. If $fh=fk$, the two maps $h,k$ define parallel maps
in the slice over the codomain of $f$, and must coincide.
Thus $f$ is monic. In particular finite endomorphism monoids
are groups, but the preceding family shows that their orders
remain unbounded.
\end{proof}

\begin{claim}{C.5}{Proved}{A direct verification of the perfect path replacement}
For each small $S$ there is a strictly functorial coaugmented
replacement $R_S:\operatorname{Cat}/S\to\operatorname{Cat}/S$
whose values are $W_\infty$-smooth and proper and whose
coaugmentation has underlying map in $W_\infty$.
\end{claim}
\begin{proof}
Let $J$ be the integer fence, with an arrow from each even integer
to each adjacent odd integer. Let $\mathcal P(S)$ have the
eventually constant functors $J\to S$ and all natural transformations.
The left and right endpoint functors are $\sigma,\beta$. These
are defined on morphisms because naturality on identity tails
forces their components to be eventually constant too.

We check the incoming and outgoing lifting categories of
$(\sigma,\beta):\mathcal P(S)\to S\times S$ directly. For an
outgoing lift from $\gamma$ with prescribed endpoint arrows
$\ell\to\ell'$ and $r\to r'$, choose an odd left boundary $a$
and an even right boundary $b$, far enough out that $\gamma$
is already constant there. Use $\gamma$ on $[a+1,b]$, use
$\ell'$ on the left tail and $r'$ on the right tail, and use
the prescribed arrows at the joins $a+1\to a$ and $b\to b+1$.
The transformation from $\gamma$ is the identity on the middle
interval and the prescribed maps on the tails. In the full
stage of lifts stationary outside these boundaries this object
is initial: the map to another lift is its given transformation
on the middle interval and the identity on the two tails.
Naturality at the joins follows from the other lift's two
naturality squares; uniqueness follows from the identity middle
components of the universal lift. For incoming lifts choose
an even left boundary and an odd right boundary and reverse
this construction. It supplies a terminal object in each stage.
Every finite simplex of either lifting category occurs in one
stage, and enlarging the boundaries is filtered. Filtered
colimits of simplicial sets preserve weak equivalences.
Both full lifting categories are therefore weakly contractible.
The lifting-category criterion for smoothness and properness
proves the asserted two properties of the endpoint functor
\citep[Theorem~4.4.36(ii'), and its opposite]{RESCisinski2019}.

For $f:A\to S$ set
\[
 R_S(f)=\bigl(A\times_{S,\sigma}\mathcal P(S)
           \xrightarrow{\beta}S\bigr),
 \qquad \eta_f(a)=(a,\operatorname{const}_{f(a)}).
\]
The map to $A\times S$ is a base change of the endpoint functor,
and the projection $A\times S\to S$ is smooth and proper.
These properties are stable under composition and base change,
so $R_S(f)$ is perfect.

There is also a direct proof that the projection $R_S(f)\to A$
is weak. Restrict to a finite fence interval containing the
support of a path and extend constantly outside it. Deleting
its rightmost leaf, and then extending constantly from the
neighbor, gives a retraction of finite diagram categories.
The edge to the deleted leaf gives a natural transformation
in one of the two directions between this composite and the
identity. Repeat to contract the finite fence to its leftmost
vertex. Throughout, that vertex is $f(a)$ and the transformations
are the identity on $a$. Thus the finite-stage projection to
$A$ is a category-homotopy equivalence, with the constant-path
section. Passing to the filtered union proves the full
projection weak. Its section $\eta_f$ is weak by two-out-of-three.
For $u:A\to A'$ over $S$, send $(a,\gamma)$ to $(u(a),\gamma)$.
These formulas preserve composition and the coaugmentation
strictly. This supplies the promised functoriality and completes
an independent verification of the portion of BISITE B.19 used here.
\end{proof}

\begin{claim}{C.6}{Proved}{All intermediate perfect subcategories compute the same localization}
Let $W$ be any fundamental localizer. For either
$\mathcal P_S=\operatorname{Parf}_{W_\infty}S$ or
$\mathcal P_S=\operatorname{Parf}_W S$, and every full
$\mathcal P_S\subseteq\mathcal H\subseteq\operatorname{Cat}/S$,
both inclusions induce equivalences after localization at $W_S$,
ordinarily and as $\infty$-categories. This applies in particular
to all $W_n$ and every homology localizer.
\end{claim}
\begin{proof}
The minimal-localizer theorem gives $W_\infty\subseteq W$
\citep[Theorem~6.1.18]{RESCisinski2006}. The lifting tests also
give $\operatorname{Parf}_{W_\infty}\subseteq\operatorname{Parf}_W$.
C.5 therefore provides a functor $R_S$ taking every object into
either chosen full subcategory, and a natural transformation
$\eta:1\to R_S$ with components in $W_S$. Naturality and
two-out-of-three show that $R_S$ preserves $W_S$.
For an object of $\mathcal H$, both ends of $\eta$ lie in
$\mathcal H$, so the transformation is internal to it by
fullness. The same holds in $\mathcal P_S$.
After localization, these transformations are natural
equivalences and make $R_S$ an inverse to each inclusion.
No assertion that all perfect maps are intrinsic for $W$
enters this argument. In particular the failure of that
intrinsic assertion for homology localizers does not obstruct
the present conclusion.
\end{proof}

\begin{claim}{C.7}{Proved}{The precise conclusion at XII ms.~240--241}
The left factorization hypothesis at XII ms.~240 implies the
three fibre axioms as printed at ms.~215--218. For arbitrary
$g:T\to S$, pullback on intrinsic fibrant objects defines a
derived reindexing functor on the localized slices. For a
strict pullback $Y=X_1\times_S X_2$ with $X_1\to S$ intrinsic,
both displayed relative homotopy-fibre identities at ms.~241
hold. These assertions do not require every terminal map to
be intrinsic. The weaker existence condition in footnote~128,
that intrinsic maps be pullbackable, suffices.
\end{claim}
\begin{proof}
C.1 and C.2 prove the three axioms. Pullback along $g$ takes
intrinsic fibrations to intrinsic fibrations and preserves
their weak arrows by C.1, so it descends to their localizations.
C.2 identifies these with the localized full slices. The
canonical isomorphisms between iterated pullbacks provide the
coherence for successive base changes.

The first identity is computed by $X_1\to S$ itself, which
is already fibrant. For the second, choose a factorization
$X_2\xrightarrow w\bar X_2\to S$ with $w\in W$ and the
second map intrinsic. The comparison
$X_1\times_SX_2\to X_1\times_S\bar X_2$ belongs to $W$
by the definition of the intrinsic map $X_1\to S$.
It therefore identifies the raw pullback with the replacement
calculation in the localized slice over $X_1$, exactly as
required. This also proves independence of replacement
through the equivalences in C.2.

Inspecting these proofs, every pullback used is along an
intrinsic map, or along a composite or a base change of one.
Those maps need only be pullbackable; arbitrary pullbacks
elsewhere in $\mathcal M$ are unnecessary. The terminal
object needed in the localization argument is $1_S$ in the
slice. Thus the caution in footnote~130 about all objects
being globally fibrant is unnecessary for the stated
relative assertions. The word ``cofibrations'' in
footnote~129 has the opposite variance from the displayed
factorization and should read ``fibrations.'' The two
mislabelled arrows at ms.~241 are $\bar p_2$, not $p_2$
and $\bar q_2$, respectively, as follows from their domains
and codomains. These are mathematical typing corrections;
no additional palaeographical claim is made.
\end{proof}

\begin{claim}{C.8}{Open}{The remaining minimum problem}
The least number of objects of a finite category with pullbacks
admitting (a) and (b) but not (c) has not been determined here.
C.3 proves an upper bound of five. The unbounded-arrow
classification required to exclude all smaller finite
categories has not been supplied. No exhaustive enumeration
through six objects is claimed.
\end{claim}
